% P11 paper module: R28 fixed Gamma crossblock and constraint-normal mismatch.

\subsection{Fixed Gamma crossblock and constraint-normal mismatch}
\label{subsec:o3aa-fixed-gamma-crossblock}

R27 removes the terminal parameter from the inverse-functional-calculus branch.  At a
fixed baseline $0<R<S<T_0$, the remaining limit geometry has two distinct components:
a tangential constrained-Gamma offblock on the effective higher-jet hyperplane and a
one-dimensional mismatch of the Riesz normals of the hard constraint
$\beta^{(0)}=0$.

For $X\in\{R,S\}$ put
\[
 B_X:=G_{X,T_0}^-
\]
and work in the corresponding baseline-whitened source Hilbert space.  Thus
\[
 W:=B_S^{1/2}J_{R,S}B_R^{-1/2}
\]
is an isometry.  Define
\[
 \widehat\beta_X(x):=\beta_X^{(0)}(B_X^{-1/2}x)
\]
and let $r_X$ be its Riesz vector:
\[
 \widehat\beta_X(x)=\langle x,r_X\rangle_X.
\]
Then the R27 effective hyperplane is
\[
 V_X=r_X^\perp.
\tag{O3AA.1}\label{eq:o3aa-VX}
\]
Pullback compatibility of the first boundary jet gives
\[
 \boxed{W^*r_S=r_R.}
\tag{O3AA.2}\label{eq:o3aa-normal-pullback}
\]

Put
\[
 N:=\Ran W,
 \qquad
 G:=N^\perp,
 \qquad
 s:=s_{R,S,T_0}:=(I-WW^*)r_S.
\tag{O3AA.3}\label{eq:o3aa-s}
\]
Then
\[
 \boxed{r_S=Wr_R+s,\qquad Wr_R\perp s,}
\tag{O3AA.4}\label{eq:o3aa-normal-split}
\]
and therefore
\[
 \|r_S\|^2=\|r_R\|^2+\|s\|^2.
\tag{O3AA.5}\label{eq:o3aa-normal-norm}
\]

\begin{proposition}[Graph form of the constrained complement]
\label{prop:o3aa-graph-complement}
Let
\[
 M:=WV_R,
 \qquad
 K:=V_S\cap M^\perp.
\]
Then
\[
 \boxed{
 K=
 \left\{
 Tg:=g-
 \frac{\langle g,s\rangle}{\|r_R\|^2}Wr_R
 :g\in G
 \right\}.
 }
\tag{O3AA.6}\label{eq:o3aa-K-graph}
\]
The map $T:G\to K$ is a bounded isomorphism whose inverse is the orthogonal projection
$P_G|_K$.
\end{proposition}

\begin{proof}
Since $N=WV_R\oplus\mathbb C Wr_R$ orthogonally,
\[
 M^\perp=\mathbb C Wr_R\oplus G.
\]
Thus $k\in K$ has a unique representation $k=aWr_R+g$ with $g\in G$.  The condition
$k\perp r_S$, together with~\eqref{eq:o3aa-normal-split}, gives
\[
 a\|r_R\|^2+\langle g,s\rangle=0,
\]
which is exactly~\eqref{eq:o3aa-K-graph}.  The inverse statement is immediate from the
orthogonal decomposition.
\end{proof}

The graph description converts R27-F into a prescribed rank-one test on the ordinary
baseline complement $G$.

Let $\Lambda_X$ be the full baseline-whitened Gamma Riesz operator:
\[
 \langle\Lambda_Xx,y\rangle
 =\widetilde{\mathfrak c}_{\Gamma,X}[x,y].
\]
Since $q_{T_0}^X=\mathfrak c_{\Gamma}+\sigma_{T_0}$ with $\sigma_{T_0}\ge0$,
\[
 0<\Lambda_X\le I,
 \qquad
 W^*\Lambda_SW=\Lambda_R.
\tag{O3AA.7}\label{eq:o3aa-Lambda}
\]
Define
\[
 \mathcal A_\Gamma
 :=P_{V_R}W^*\Lambda_S|_G:G\to V_R,
 \qquad
 a_\Gamma:=P_{V_R}\Lambda_Rr_R.
\tag{O3AA.8}\label{eq:o3aa-Aa}
\]

\begin{proposition}[Prescribed rank-one criterion]
\label{prop:o3aa-rankone}
With the convention that the Hilbert inner product is linear in the first argument,
\[
 \boxed{
 \widetilde{\mathfrak c}_{\Gamma,S}[Wf,Tg]
 =\langle f,\mathcal A_\Gamma g\rangle
 -\frac{\langle s,g\rangle}{\|r_R\|^2}
  \langle f,a_\Gamma\rangle
 }
\tag{O3AA.9}\label{eq:o3aa-cross-form}
\]
for $f\in V_R$, $g\in G$.  Hence
\[
 \boxed{
 \mathscr C_\Gamma^{R,S,T_0}=0
 \quad\Longleftrightarrow\quad
 \mathcal A_\Gamma g
 =\frac{\langle g,s\rangle}{\|r_R\|^2}a_\Gamma
 \quad(g\in G).
 }
\tag{O3AA.10}\label{eq:o3aa-rankone-criterion}
\]
In particular, if $s=0$, then
\[
 \mathscr C_\Gamma=0\iff\mathcal A_\Gamma=0.
\tag{O3AA.11}\label{eq:o3aa-rankone-szero}
\]
\end{proposition}

\begin{proof}
Insert~\eqref{eq:o3aa-K-graph} in the Gamma form.  The first term is represented by
$\mathcal A_\Gamma$.  Gamma pullback compatibility gives
\[
 \widetilde{\mathfrak c}_{\Gamma,S}[Wf,Wr_R]
 =\widetilde{\mathfrak c}_{\Gamma,R}[f,r_R]
 =\langle f,a_\Gamma\rangle.
\]
Because the second slot is conjugate-linear, the coefficient
$\langle g,s\rangle$ in~\eqref{eq:o3aa-K-graph} appears as $\langle s,g\rangle$ in
\eqref{eq:o3aa-cross-form}.  Vanishing for all $f$ is therefore equivalent to
\eqref{eq:o3aa-rankone-criterion}.
\end{proof}

Now let $L_X$ be the constrained Gamma operator on $V_X$ represented by the R27 limit:
\[
 \langle L_Xx,y\rangle
 =\widetilde{\mathfrak c}_{\Gamma,X}[x,y].
\]
By Corollary~\ref{cor:o3z-whitened-mosco},
\[
 a_0I\le L_X\le I.
\tag{O3AA.12}\label{eq:o3aa-L-bounds}
\]
With respect to $V_S=M\oplus K$,
\[
 \boxed{
 L_S=
 \begin{pmatrix}
 WL_RW^*&\mathscr C_\Gamma^*\\
 \mathscr C_\Gamma&D
 \end{pmatrix}.
 }
\tag{O3AA.13}\label{eq:o3aa-L-block}
\]
R27 also identifies
\[
 T_{X,\infty}=L_X^{-1/2}P_{V_X}
\tag{O3AA.14}\label{eq:o3aa-T-limit}
\]
and
\[
 D_\infty^-:=T_{S,\infty}W-WT_{R,\infty}.
\tag{O3AA.15}\label{eq:o3aa-Dinf}
\]

\begin{theorem}[Tangential inverse-root equivalence]
\label{thm:o3aa-tangential-equivalence}
On the effective higher-jet domain,
\[
 \boxed{
 D_\infty^-|_{V_R}=0
 \quad\Longleftrightarrow\quad
 \mathscr C_\Gamma^{R,S,T_0}=0.
 }
\tag{O3AA.16}\label{eq:o3aa-tangential-equivalence}
\]
\end{theorem}

\begin{proof}
If $\mathscr C_\Gamma=0$, then~\eqref{eq:o3aa-L-block} is block diagonal and
$L_SW=WL_R$ on $V_R$.  Continuous functional calculus gives
\[
 L_S^{-1/2}W=WL_R^{-1/2},
\]
which is the forward implication.

Conversely assume the inverse-root intertwining on $V_R$.  Apply it to
$L_R^{1/2}x$ and multiply by $L_S^{1/2}$ to get
\[
 WL_R^{1/2}x=L_S^{1/2}Wx.
\]
Applying $L_S^{1/2}$ once more yields
\[
 L_SWx=WL_Rx.
\]
Thus $M=WV_R$ is invariant under the selfadjoint operator $L_S$, hence reducing, and
the offblock in~\eqref{eq:o3aa-L-block} vanishes.
\end{proof}

Theorem~\ref{thm:o3aa-tangential-equivalence} supplies the converse that was not part of
Proposition~\ref{prop:o3z-gamma-block}: vanishing of the limiting inverse square-root
intertwining defect on the higher-jet domain is exactly equivalent to the fixed Gamma
offblock being zero.

\begin{theorem}[Full-space decomposition and normal-mismatch lower bound]
\label{thm:o3aa-full-equivalence}
Let
\[
 e_R:=r_R/\|r_R\|.
\]
Then
\[
 \boxed{
 \|D_\infty^-e_R\|
 \ge\frac{\|s\|}{\|r_S\|}.
 }
\tag{O3AA.17}\label{eq:o3aa-normal-lower}
\]
Consequently
\[
 \boxed{
 D_\infty^-=0
 \text{ on the full standardized source space}
 \quad\Longleftrightarrow\quad
 \mathscr C_\Gamma=0
 \text{ and }s=0.
 }
\tag{O3AA.18}\label{eq:o3aa-full-equivalence}
\]
\end{theorem}

\begin{proof}
Because $e_R\in V_R^\perp$, one has $T_{R,\infty}e_R=0$, hence
\[
 D_\infty^-e_R=L_S^{-1/2}P_{V_S}We_R.
\]
The upper bound $L_S\le I$ implies $L_S^{-1/2}\ge I$.  Moreover
\[
 \left|
 \left\langle We_R,\frac{r_S}{\|r_S\|}\right\rangle
 \right|
 =\frac{\|r_R\|}{\|r_S\|}
\]
by~\eqref{eq:o3aa-normal-pullback}.  Therefore
\[
 \|P_{V_S}We_R\|^2
 =1-\frac{\|r_R\|^2}{\|r_S\|^2}
 =\frac{\|s\|^2}{\|r_S\|^2},
\]
which proves~\eqref{eq:o3aa-normal-lower}.  Thus the normal direction intertwines iff
$s=0$.  Combining this with Theorem~\ref{thm:o3aa-tangential-equivalence} and the
orthogonal decomposition $V_R\oplus\mathbb Ce_R$ proves
\eqref{eq:o3aa-full-equivalence}.
\end{proof}

\begin{remark}[R28 firewall]
\label{rem:o3aa-firewall}
R28 shows that the fixed residue of the R27 inverse-root branch has two layers:
\begin{enumerate}
 \item the tangential Gamma offblock $\mathscr C_\Gamma$;
 \item the constraint-normal mismatch $s=(I-WW^*)r_S$.
\end{enumerate}
Either nonvanishing gives a persistent inverse-functional-calculus defect.  This is
still modulus-layer information.  The R14 countermodel continues to forbid promotion
to the relative polar gauge without separate control of the polar factors.
\end{remark}

\begin{openproblem}[R28-F: fixed tangential and normal tests]
\label{open:o3aa-fixed-tests}
Decide, for the concrete fixed P11 baseline,
\[
 \boxed{
 \mathcal A_\Gamma
 \stackrel{?}{=}
 \frac{a_\Gamma\otimes s^*}{\|r_R\|^2}
 }
\]
and
\[
 \boxed{s=(I-WW^*)r_S\stackrel{?}{=}0.}
\]
The first question is the tangential constrained-Gamma invariance problem; the second
asks whether the hard first-boundary constraint has compatible baseline Riesz normals
under source enlargement.  No terminal parameter remains in either question.
\end{openproblem}
